% Abstract (lead, 23 Nov 2026). Every number is copied from ch01_intro.tex (traced in ch01_intro_NOTES.md) or RELATED_WORK.md l88 / l92.
Electroluminescence (EL) images show cracks and inactive regions in solar cells that the eye cannot see, and many papers report
segmentation scores for them. Those scores do not compare: of 22 papers with pixel-level EL numbers read in full, none states a
module-disjoint split and 18 give no interval, seed count or repeat. This thesis builds a benchmark for 24-class pixel segmentation
of EL cells and uses it to test what actually improves crack segmentation.

A dataset audit shows that strict crack IoU measures the annotators' line-width convention: the label release alone moves the
baseline's strict crack IoU by +0.045, more than any method gain over it on the pool. The protocol therefore splits the 532 pool
cells by physical module, decides on crack relF1@2 with a 2~px tolerance and a margin of 0.010, reports seed-aware intervals, and
reads each sealed set once under a rule registered before any prediction. Under this protocol, two models from the EL literature
fall below a U-Net++ baseline, and no trained arm read at five seeds was confirmed under its registered rule.

On the sealed set E1, the final recipe (HRNet-W18 U-Net++ with SWAD weight averaging and D4 test-time augmentation) beats the
baseline on crack relF1@2 by +0.0430 [+0.0013, +0.1043], confirmed with all guards passing; the recipe is averaged over seeds 5 to 9,
and 101 of the 106 cells come from one source. A registered few-shot test inside E1 shows that labelled cells of the target layouts
help: the recipe trained with them beats its zero-shot use by +0.2213 [+0.0977, +0.3504], carried by the half-cut cells. The recipe costs about
0.17~s of network time per cell on one GPU. The benchmark protocol, split lists and scoring code are released; images, masks and
weights are not.
